\section*{Capítulo 5}

\solucion{5.1} $F_0=50\,e^{0.04\times0.5}=50\times1.020201=51.01$.

\solucion{5.2} El precio justo es $100e^{0.05}=105.13>102$: el forward está barato. Estrategia inversa (\emph{reverse cash and carry}): hoy se vende en corto el activo (+100), se invierten los 100 al 5\,\% y se compra forward a 102 (costo cero). En $T$: la inversión vale 105.13; se pagan 102 por el activo comprado a plazo y con él se cierra la venta en corto. Ganancia: $105.13-102=3.13$ sin riesgo.

\solucion{5.3} $F=3.40\times1.045/1.040=3.4163$. El dólar cotiza con \emph{prima} (más caro a plazo) porque la tasa en dólares es menor que la tasa en soles.

\solucion{5.4} Liquidación para el comprador de soles: $100/3.30-100/3.50=30.303-28.571=+1.732$ millones de dólares. El sol se apreció por encima de lo pactado: el banco paga a la empresa US\$ 1.73 millones.

\solucion{5.5} Para la parte corta: $f=(89.38-95.00)e^{-0.0412\times0.25}=-5.62\times0.98975=-5.56$ dólares por barril. Es un pasivo para PETROPERÚ, compensado por el mayor valor del inventario.

\solucion{5.6} El forward no es un pronóstico: está determinado por la paridad cubierta, es decir, por el diferencial de tasas entre soles y dólares. Que cotice por debajo del spot solo refleja que la tasa en soles es menor que la de dólares. Además, el objetivo de cubrirse no es acertar la dirección del tipo de cambio sino eliminar la variabilidad del resultado.

\solucion{5.7} Porque en la operación subyacente PETROPERÚ vende dólares al banco para obtener soles, y el precio al que el banco compra dólares es el de compra (3.4367). Usar el de venta sobrestimaría los soles que la empresa obtiene por cada dólar.
